\documentclass[11pt, a4paper]{amsart}

\RequirePackage{amssymb,amsfonts, amsmath, amsthm,latexsym}
\RequirePackage[british]{babel}
\RequirePackage{version,graphicx,times,titletoc,dirtytalk}
\RequirePackage[utf8]{inputenc}
\RequirePackage{xcolor, url, hyperref}

\setlength{\oddsidemargin}{.30in}
\setlength{\evensidemargin}{.30in} \setlength{\textwidth}{5.7in}
\setlength{\parskip}{5pt} \setlength{\labelsep}{5pt}
\setlength{\parindent}{0pt} \setlength{\medskipamount}{3ex}


\newtheorem {thm}{Theorem}
\newtheorem {cor}[thm]{Corollary}
\newtheorem {lem}[thm]{Lemma}
\newtheorem {prop}[thm]{Proposition}
\newtheorem {conj}[thm]{Conjecture}
\theoremstyle{definition}
\newtheorem {defi}[thm]{Definition}
\newtheorem {rem}[thm]{Remark}
\newtheorem {exa}[thm]{Example}

\DeclareMathOperator{\ord}{ord}
\DeclareMathOperator{\Gal}{Gal}
\DeclareMathOperator{\Aut}{Aut}
\DeclareMathOperator{\Hom}{Hom}
\DeclareMathOperator{\Ind}{Ind}
\DeclareMathOperator{\re}{Re}
\DeclareMathOperator{\im}{Im}
\DeclareMathOperator{\rad}{rad}
\DeclareMathOperator{\Li}{Li}
\DeclareMathOperator{\id}{id}
\DeclareMathOperator{\N}{N}
\DeclareMathOperator{\lcm}{lcm}
\DeclareMathOperator{\dens}{dens}
\DeclareMathOperator{\rk}{rank}
\DeclareMathOperator{\Id}{Id}
\DeclareMathOperator{\tr}{tr}
\DeclareMathOperator{\Spec}{Spec}
\DeclareMathOperator{\Res}{Res}
\DeclareMathOperator{\radd}{rad}

\newcommand{\F}{\mathbb{F}}
\newcommand{\Q}{\mathbb{Q}}
\newcommand{\NN}{\mathbb{N}}
\newcommand{\Z}{\mathbb{Z}}
\newcommand{\PP}{\mathbb{P}}
\newcommand{\p}{\mathfrak{p}}
\newcommand{\q}{\mathfrak{q}}
\newcommand{\pord}{\ord_\p}
\newcommand{\lval}{v_\ell}
\newcommand{\rank}{\text{rank}}
\newcommand{\OO}{\mathcal{O}}
\newcommand{\End}{\mathrm{End}}
\newcommand{\tors}{\mathrm{tors}}

\newcommand{\ch}{\mathrm{char}}

\renewcommand{\geq}{\geqslant}
\renewcommand{\leq}{\leqslant}
\renewcommand{\l}{\elleqslant}

\newcommand{\group}[1]{\left\langle#1\right\rangle}

\newcommand{\zl}[2][\ell]{\zeta_{{#1}^{#2}}}



\title{Linear relations among radicals}
\author{Antonella Perucca}
\date{}

\keywords{Kneser's theorem, Kummer theory, radicals, entanglement.}
\subjclass[2000]{Primary: 12J99; Secondary: 11R18.}

\begin{document}


\maketitle

\begin{center}
     \textit{In memory of Marc Rybowicz}
\end{center}

\begin{abstract}
Let $K$ be a field, fix an algebraic closure $\overline{K}$, and let $G$ be a subgroup of $\overline{K}^\times$. We are able to give a closed formula for the ratio between the degree $[K(G):K]$ and the index $|GK^\times:K^\times|$, provided that the latter is finite. Our formula  explains all the $K$-linear relations among radicals, which (beyond the ones stemming from the multiplicative group $GK^\times/K^\times$) are generated by relations among roots of unity and single radicals. Our work builds on results by Rybowicz, which in turn are based on work by Kneser and Schinzel. 
\end{abstract}




\section{Introduction}

We let $K$ be a field, for which we fix an algebraic closure $\overline{K}$. We consider the \emph{radicals} over $K$, by which we mean the elements $\alpha\in \overline{K}^\times$ for which there exists a positive integer $n$ -- coprime to the characteristic of $K$ -- such that $\alpha^n\in K^\times$. 
We fix a group $G$ of radicals such that the index $|GK^\times :K^\times|$ is finite. We investigate the degree $[K(G):K]$ (the extension $K(G)/K$ is separable, but in general not Galois). More precisely, we compare the above degree and  index.
If all $K$-linear relations among the radicals in $G$ stem from multiplicative relations, then the above-mentioned degree and index will be the same. Else, we have a phenomenon that is called \emph{entanglement} (of radicals). We denote by $n$ the smallest positive integer such that $G^n\subseteq K^\times$ and by $z$ the product of the odd prime divisors $p$ of $n$ such that $\zeta_p\notin K^\times$ and $\zeta_p\in GK^\times$. If $H$ is a multiplicative group and $m$ is a positive integer, we write $\mu_m(H)$ for the subgroup of $H$ consisting of roots of unity of order dividing $m$. The main result of this paper is the following:

\begin{thm}\label{mainthm}
If $n$ is odd, we have
$$\frac{[K(G):K]}{|GK^\times:K^\times|}=\frac{[K(\zeta_{z}):K]}{|\mu_{n}(G K^\times)\cap K(\zeta_{z})^\times: \mu_{n}(K^\times)|}\,.$$
If $n$ is even, writing $n=2^f n'$ where $f$ is a positive integer and $n'$ is an odd integer, the ratio $\frac{[K(G):K]}{|GK^\times:K^\times|}$ equals 
$$\frac{[K(\zeta_{z}):K] \cdot 2^{-\Delta}}{|\mu_{n'}(G K^\times)\cap K(\zeta_{z})^\times: \mu_{n'}( K^\times)|\cdot |\mu_{2^{f+1}}(G K^\times)(G K^\times\cap\sqrt{K^\times})\cap K(\zeta_z)^\times: K^\times|}$$
where $\Delta$ is the non-negative integer from Definition \ref{Delta}.
\end{thm}

This result allows to completely understand the $K$-linear relations among the radicals in $GK^\times$: for $p$ an odd prime such that $\zeta_p\in GK^\times$ and $\zeta_p\notin K^\times$, calling $d_p=[K(\zeta_p):K]$, we have 
$$\zeta_p^{d_p}, \ldots, \zeta_{p}^{p-1}\in 1K+\zeta_p K+\ldots+ \zeta_p^{d_p-1} K=K(\zeta_p)\,.$$
Moreover, the fact that certain powers of $\zeta_{p^{v_p(n)}}$ of order larger than $p$ (and contained in $GK^\times$) may be contained in $K(\zeta_p)^\times$ leads to further $K$-linear relations among the roots of unity.

Beyond this phenomenon, the elements in $\mu_{n'}(G K^\times)\cap K(\zeta_{z})^\times$
are powers of $\zeta_{n'}$ and also $K$-linear combinations of powers of $\zeta_{z}$,  and equating the two expressions gives rise to a $K$-linear relation among roots of unity. Similarly, the elements in 
$\mu_{2^{f+1}}(G K^\times)(GK^\times\cap\sqrt{K^\times})\cap K(\zeta_z)^\times$ provide $K$-linear relations between single elements of $G$ whose square is in $K^\times$ and powers of $\zeta_{2^{f+1}z}$. Finally, as explained in Section \ref{sec:better2}, the term $2^{-\Delta}$ stems from $K$-linear relations among roots of unity of order dividing $2^{f+1}$, and possibly an additional relation
$$1+\zeta_{2^w}\in 1K+\zeta_4K$$
where $w$ is largest integer such that $\zeta_{2^w}+\zeta_{2^w}^{-1}\in K^\times$ (provided that such largest integer exists).

From the main result we deduce the following general property:

\begin{thm}\label{divisibility}
The degree $[K(G):K]$ divides 
$$\frac{1}{z} \cdot [K(\zeta_z):K]  \cdot |GK^{\times}:K^{\times}|\,.$$
\end{thm}

We conclude the paper by investigating the growth of radical extensions (see Theorems \ref{eventual} and \ref{MAMA}). 

We suppose that $n$ is coprime to the characteristic of $K$ but this assumption is not necessary. Indeed, if $\ch(K)=p$, the extension $K(G)/K(G^{p^{v_p(n)}})$ is purely inseparable and with degree the $p$-part of $|GK^\times: K|$ (see \cite[Corollary 9.2, Chapter VI]{LangAlgebra}). We deduce that
$$\frac{[K(G):K]}{|GK^\times:K^\times|}=\frac{[K(G^{p^{v_p(n)}}):K]}{|G^{p^{v_p(n)}}K^\times:K^\times|}$$
so we have reduced to the case where $n$ is coprime to $p$.

We also remark that to study the $K$-linear relations of (finitely many) radicals we may consider the group $G$ that they generate, so our assumption that $|GK^\times:K^\times|$ is finite is not restrictive.

Our results build on two theorems by Rybowicz (Theorems \ref{thmKR} and \ref{Ryb2}) which express the degree $[K(G):K]$ also in the cases in which the famous theorem by Kneser (Theorem \ref{Kneserthm}) cannot be applied. Other results that we make use of is a lemma by Schinzel about radicals to extend the base field to $K(\zeta_4)$ (Lemma \ref{lemminoSchinzel}) and Schinzel's Theorem on abelian radical extensions (Theorem \ref{Schinzel-abelian}).
There is a vast literature on radical extensions however our general result seem to be new. We mention for example also \cite{Koch2} by 
Halter-Koch, \cite{BarreraVelez} by Barrera Mora and Vélez, and  \cite{Albu} by Albu. In \cite{Lenstra} Lenstra investigated entanglements introducing the entanglement group: this group was studied also by Palenstijn \cite{Palenstijn} and by the author with Sgobba and Tronto \cite{MAMA}. The entanglement has also been studied by Lenstra, Moree and Stevenhagen in \cite{LMS}. Recently, the author with Chan, Pajaziti, and Perissinotto established further results on the entanglement, see \cite{EntRad}.

\subsection*{Acknowledgments} We thank Daniel Gil-Munoz for discussions which also lead to a special case of the main theorem, Zeev Rudnick for a general discussion about radical extensions, and Fritz Hörmann for many useful comments.

\section{Kneser's Theorem and Kummer theory}\label{sec:KneserKummer}

We let $K$ be a field and fix an algebraic closure $\overline{K}$. We  let $G$ be a group of radicals over $K$ such that the index $|G K^\times :K^\times|$ is finite and coprime to the characteristic of $K$. We let $n$ be a positive integer, which we  suppose to be minimal, such that $G^n$ is contained in $K^\times$ (thus, $n$ is coprime to the characteristic of $K$). 

For any positive integer $m$ that is coprime to $\ch(K)$ we fix some root of unity $\zeta_m$ in $\overline{K}^\times$ of order $m$ (with a coherent choice, namely that if $m,M$ are positive integers such that $m\mid M$ then we have $\zeta_M^{M/m}=\zeta_m$). 
If $H$ is a subgroup of $\overline{K}^\times$ and $m$ is a positive integer, we write $\mu_m(H)$ for the group of roots of unity in $H$ whose order divides $m$.

\begin{rem}\label{remexact}
The following sequence, induced by the exponentiation by $n$, is exact:
\begin{equation}\label{exact}
1 \rightarrow \mu_n(GK^\times)K^\times/K^\times \rightarrow GK^\times/K^\times \rightarrow G^nK^{\times n}/K^{\times n} \rightarrow 1\,.
\end{equation}
Indeed, if $(ga)^n=b^n$ for some $g\in G$ and for some $a,b\in K^\times$, then $g^n\in K^{\times n}$ and hence $g\in \mu_n(GK^\times)K^\times$.
Moreover, the following sequence is exact 
\begin{equation*}\label{exactroot}
1 \rightarrow \mu_n(K^\times) \rightarrow \mu_n(GK^\times) \rightarrow \mu_n(GK^\times)K^\times/K^\times \rightarrow 1
\end{equation*}
because we have $\mu_n(K^\times)=\mu_n(GK^\times)\cap K^\times$.
We deduce that 
\begin{equation}\label{exactall}
|GK^\times:K^\times|=|G^nK^{\times n}:K^{\times n}|\cdot | \mu_n(GK^\times): \mu_n(K^\times)|\,.
\end{equation}
\end{rem}

We rely on the famous result by Martin Kneser from \cite{Kneser}:

\begin{thm}[Kneser's Theorem]\label{Kneserthm}
We have 
$$[K(G):K]=|GK^\times:K^\times|$$
if the following two conditions hold: for every odd prime $p$ we have $\zeta_p\in K^\times$ or $\zeta_p\notin GK^\times$; we have $\zeta_4\in K^\times$ or $1\pm \zeta_4\notin G K^\times$. 
\end{thm}

Our main result implies the following: 

\begin{prop}
The two conditions in Kneser's theorem are necessary. 
\end{prop}
\begin{proof}
We rely on Theorem \ref{mainthm}.
If $z>1$, then $[K(\zeta_z):K]\leq \varphi(z)$ while 
$|\mu_{n'}(G K^\times)\cap K(\zeta_{z})^\times:K^\times|$ is divisible by $z$ hence there is some prime $p\mid z$ such that the $p$-adic valuation of $[K(G):K]/|GK^\times:K^\times|$ is non-zero. Now suppose that $z=1$ but that the second condition in Kneser's theorem does not hold: we prove that the $2$-adic valuation of $[K(G):K]/|GK^\times:K^\times|$ is non-zero. To study this $2$-adic valuation, we may replace $G$ by $G^{n'}$ (hence $n$ becomes $2^f$). Then by Kneser's theorem over $L=K(\zeta_4)$ we have 
$[L(G):L]=|GL^\times:L^\times|$. We may conclude because $[K(G):K]/[L(G):L]$ divides $2$ while $|GK^\times:K^\times|/|GL^\times:L^\times|$ is a multiple of $4$ because (with the appropriate sign choice) the class of $(1\pm \zeta_4)\in L^\times$ has order $4$ in 
$GK^\times/K^\times$.
\end{proof}

\begin{rem}\label{Knesercondi}
If $\zeta_p\in GK^\times$ and $\zeta_p\notin K^\times$, then we have $p\mid n$. If (for a sign choice) $1\pm \zeta_4\in GK^\times$ and $\zeta_4\notin K^\times$, then $4\mid n$ because  
$(1\pm \zeta_4)^2=\pm 2\zeta_4$ and $(1\pm \zeta_4)^4=-4$  hence the order of $1\pm \zeta_4$ in $GK^\times/K^\times$ is $4$. 
\end{rem}

By the above remark, the two conditions in Kneser's theorem are satisfied if $\zeta_n\in K$.
In this case, the extension $K(G)/K$ is a \emph{Kummer extension} and its Galois group is abelian of exponent dividing $n$:

\begin{thm}[Kummer theory]
    If $\zeta_n\in K$, the groups $\Gal(K(G)/K)$ and $GK^\times/K^\times$ and $G^nK^{\times n}/K^{\times n}$ are isomorphic. In particular, we have
    \begin{equation}\label{all} [K(G):K]=|GK^\times:K^\times|=|G^nK^{\times n}:K^{\times n}|\,.
    \end{equation}
\end{thm}
\begin{proof}
The isomorphism between the former and latter group is one of the main results in Kummer theory (see \cite[Theorem 8.1, Chapter VI]{LangAlgebra}). The isomorphism between the second and the third group is a consequence of \eqref{exact} because $\mu_n(GK^\times)=\mu_n(K^\times)=\langle \zeta_n\rangle$.
\end{proof}

If $K$ and $G$ consist of real numbers, then $\mu_n(GK^\times)=\mu_n(K^\times)=\{\pm 1\}$ and the two conditions in Kneser's theorem are satisfied hence \eqref{all} holds (see also  \cite[Theorem 2.2]{Rybowicz}).

\section{The case where $n$ is an odd prime power}

We suppose that $n$ is the power of an odd prime number $p$ (thus, the characteristic of $K$ is different from $p$). 
We rely on the following result, which combines Theorem \ref{Kneserthm} (in view of  Remark \ref{Knesercondi}) and \cite[Theorem 2.3]{Rybowicz}: 

\begin{thm}[Kneser - Rybowicz]\label{thmKR}
We have
\begin{align*}
[K(G):K] &= |G^nK^{\times n}:K^{\times n}| \cdot [K(\mu_{n}(GK^\times)):K]\,.
\end{align*}
Moreover, if $\zeta_p\notin GK^\times$ or $\zeta_p\in K^\times$, then we have 
$$[K(G):K]= |GK^{\times}:K^{\times}|\,.$$
\end{thm}

We deduce that, if $\zeta_p\notin GK^\times$ or $\zeta_p\in K^\times$, then the degree $[K(G):K]$ is a power of $p$ while in the remaining case it is a power of $p$ times $[K(\zeta_p):K]$.


\begin{cor}\label{corbetterp}
If $H:=\mu_{n}(GK^\times)$, then we have
$$\frac{[K(G):K]}
{|GK^{\times}:K^{\times}|}=\frac{[K(H):K]}
{|HK^{\times}:K^{\times}|}\,.$$
If $\zeta_p\notin K^\times$ and $\zeta_p\in GK^\times$, we have  $H=\langle \zeta_{p^m}\rangle$ for some positive integer $m$. Let $m_0$ be the largest positive integer such that $\zeta_{p^{m_0}}\in K(\zeta_p)^\times$ (or $\infty$, if no such largest integer exists). Then we have
$$\frac{[K(G):K]}
{|GK^{\times}:K^{\times}|}=\begin{cases} 1 & \text{if $\zeta_p\in K^\times$ or $\zeta_p\notin GK^\times$}\\
[K(\zeta_p):K] \cdot p^{-{\min(m_0,m)}} & \text{otherwise}\,.
\end{cases}
$$
\end{cor}
\begin{proof}
    Combining Theorem \ref{thmKR} and  \eqref{exactall} we get
$$\frac{[K(G):K]}
{|GK^{\times}:K^{\times}|}=\frac{[K(\mu_{n}(GK^\times)):K]}{| \mu_{n}(GK^\times): \mu_{n}(K^\times)|}\,.$$
From $H\cap K^{\times}=\mu_{n}(K^\times)$ we deduce that $|H: \mu_{n}(K^\times)|=|HK^{\times}:K^{\times}|$.
By Theorem \ref{thmKR} we are left to deal with the case $\zeta_p\notin K^\times$ and $\zeta_p\in GK^\times$. We may conclude because $|H: \mu_{n}(K^\times)|=p^m$ while $[K(H):K]=[K(\zeta_p):K]p^{\max(m-m_0,0)}$.
\end{proof}

All entanglement stems from $K$-linear relations among roots of unity. The $K$-linear relation
$$1+\zeta_p+\zeta_p^2+\cdots +\zeta_p^{p-1}=0$$
has to be counted only if $\zeta_p\notin K^\times$, and it is only relevant if $\zeta_p\in GK^\times$.
If $d_p:=[K(\zeta_p):K]$, the minimal polynomial of $\zeta_p$ gives a $K$-linear relation among $1,\zeta_p,\ldots,\zeta_p^{d_p}$, and all the powers $\zeta_p^i$ for $i\geq d_p$ are $K$-linear combinations of the roots of unity $1,\zeta_p,\ldots, \zeta_p^{d_p-1}$.
We also have 
$$(\zeta_{p^m})^{p^{j+\max(m-m_0,0)}}\in 1K+\zeta_pK+\cdots +\zeta_p^{d_p-1}K \qquad \text{for}\quad j\geq 0$$
because any element in $K(\zeta_p)$ is of this form.
These $K$-linear relations generate all others for $K(G)$ (beyond those stemming from the group $GK^\times/K^\times$) because by the above result they already explain the degree of $K(G)/K$.

\section{The case where $n$ is a power of $2$}\label{sec:better2}

Let $n=2^f$ for some positive integer $f$. \emph{We suppose that $f\geq 2$ and $\zeta_4\notin K^\times$} (else, we already know that $[K(G):K]=|GK^\times:K^\times|$ by Theorem \ref{Kneserthm} and Remark \ref{Knesercondi}).
For every positive integer $t$ we write $\xi_{2^t}=\zeta_{2^t}+\zeta_{2^t}^{-1}$. Moreover, we let $w$ be the largest integer such that $\xi_{2^w}\in K^\times$, or set $w=\infty$ if no such largest integer exists.

The following lemma is due to Schinzel, and it also holds for $f=1$ (see \cite[Lemma 2]{Schinzel-ab}):

\begin{lem}[Schinzel]\label{lemminoSchinzel}
 The kernel of the map
$$K^{\times}/K^{\times n}\rightarrow K^{\times}K(\zeta_4)^{\times n}/K(\zeta_4)^{\times n}
$$
induced by the inclusion is generated by the class of the following element:
$$a=\begin{cases} -1 & \text{if $w>f$}\\
-\xi_{2^{w+1}}^n & \text{if $w=f$} \\
\xi_{2^{w+1}}^n & \text{if $w<f$}\,. \\
\end{cases}$$
\end{lem}

Notice that $\xi_{2^{w+1}}^n=(\xi_{2^w}+2)^{\frac{n}{2}}$. Since $\xi_{2^w}\in K^{\times}$, we always have $a^2\in K^{\times n}$. However, the class of $a$ modulo $K^{\times n}$ may have order $1$ or $2$: 

\begin{lem}\label{orders}
With the notation of Lemma \ref{lemminoSchinzel}, we have $a\notin K^{\times n}$ if and only if 
$w\geq f$ or $\ch(K)=0$ and $K\cap \Q(\zeta_{2^\infty})$ is totally real.
\end{lem}
\begin{proof}
We have $(-1)\notin K^{\times n}$ because $\zeta_4\notin K^\times$, so now suppose that $w\leq f$.
The condition $a\in K^{\times n}$ means that $(\xi_{2^{w}}+2)\gamma \in K^{\times 2}$, where $\gamma$ is a root of unity of order dividing $n/2$ for $w<f$ and of order $n$ for $w=f$. Since $\gamma\in K^\times$, we must have $\gamma\in \{\pm 1\}$. We cannot have $\gamma=1$ because $\xi_{2^{w+1}}\notin K$, so now suppose that $\gamma=-1$. 
The condition $-(\xi_{2^{w}}+2)=-\xi_{2^{w+1}}^2 \in K^{\times 2}$ holds in odd characteristic because for finite fields the product of two non-squares is a square (by Kummer theory, as $F(\sqrt{b})=F(\zeta_4)$ holds if $F$ is a finite field and  $b\in F^\times\setminus F^{\times  2}$). In characteristic $0$, the square roots of $-(\xi_{2^{w}}+2)$ are in $\Q(\zeta_{2^\infty})$ and not totally real. They cannot be in $K^\times$ if $K\cap \Q(\zeta_{2^\infty})$ is totally real. In the remaining case, $\zeta_4$ and $\xi_{2^{w+1}}$ generate the same quadratic extension of $K\cap \Q(\zeta_{2^\infty})$ hence of $K$, so by Kummer theory $-(\xi_{2^{w}}+2)$ is a square in $K^\times$.
\end{proof}

We rely on the following result, which is \cite[Theorem 2.4]{Rybowicz}, restated thanks to Lemma \ref{orders}:

\begin{thm}[Rybowicz]\label{Ryb2}
We have 
$$\frac{[K(G):K]}{|G^nK^{\times n}:K^{\times n}|}=\delta \cdot [K(\mu_n(GK^\times)):K]$$
where $\delta\in \{1,\frac{1}{2}\}$.
We have $\delta=\frac{1}{2}$ if and only if $a\in G^nK^{\times n}$ and, in the case $w<f$, additionally $1+\zeta_{2^w}\in GK^\times$ and $\ch(K)=0$ and $K\cap \Q(\zeta_{2^\infty})$ is totally real.
\end{thm}


\begin{thm}\label{better2}
Let $\delta$ be as in Theorem \ref{Ryb2} and set
$$H=\begin{cases} \mu_{2n}(GK^\times) & \text{if $w>f$}\\
\mu_{n}(GK^\times) & \text{if $w=f$ and $-\xi_{2^{w+1}}^n\notin G^nK^{\times n}$}\\
\langle 1+\zeta_{2^w}, \zeta_n\rangle & \text{if $w=f$ and $-\xi_{2^{w+1}}^n\in G^nK^{\times n}$}\\
\mu_{n}(GK^\times) & \text{if $w<f$ and $\delta=1$}\\
<1+\zeta_{2^w}>\mu_{n}(GK^\times)& \text{if $w<f$ and $\delta=1/2$}\,.\\
\end{cases}$$
Then $H$ is a subgroup of $GK^\times$ and we have
$$\frac{[K(G):K]}
{|GK^{\times}:K^{\times}|}= \frac{[K(H):K]}
{|HK^{\times}:K^{\times}|}\,.$$
\end{thm}
\begin{proof}
For $w=f$, we first prove that $H=\langle 1+\zeta_{2^w}, \zeta_n\rangle \cap GK^\times$. Observe that $\zeta_{2^{w+1}}\xi_{2^{w+1}}=1+\zeta_{2^w}$. If $-\xi_{2^{w+1}}^n\notin G^nK^{\times n}$, there is no integer $i$ such that $(1+\zeta_{2^w})\zeta_n^i\in GK^\times$ and hence $\langle 1+\zeta_{2^w}, \zeta_n\rangle \cap GK^\times=\mu_n(GK^\times)$. Else, fix $i$ such that $(1+\zeta_{2^w})\zeta_n^i\in GK^\times$. Considering that $\xi_{2^{w+1}}^2\in K^\times$, we deduce that $\zeta_n\in GK^\times$ and hence $1+\zeta_{2^w}\in GK^\times$ and we conclude.

So in all cases $H$ is a subgroup of $GK^\times$ such that $\mu_n(HK^\times)=\mu_n(GK^\times)$. We now prove that $\delta_G=\delta_{H}$.
If $w>f$, this is because $-1\in G^nK^{\times n}$ is equivalent to $\zeta_{2n}\in GK^{\times}$.
If $w=f$, this is because $-\xi_{2^{w+1}}^n\in G^nK^{\times n}$ is equivalent to $(1+\zeta_{2^w})\in \mu_n(\overline{K}^\times) GK^{\times}$ and we have $H=\langle 1+\zeta_{2^w}, \zeta_n\rangle \cap GK^\times$. If $w<f$, we observe that $\delta_G=1$ implies $\delta_{H}=1$, so suppose that $\delta_G=\frac{1}{2}$. We clearly have $1+\zeta_{2^w}\in HK^{\times}$, so we are left to prove that $\xi_{2^{w+1}}^n\in H^nK^{\times n}$. This is equivalent to $\xi_{2^{w+1}}\in <1+\zeta_{2^w},\zeta_n>K^{\times n}$ and we may conclude because $\xi_{2^{w+1}}=\zeta_{2^{w+1}}^{-1}(1+\zeta_{2^w})$ and $\zeta_{2^{w+1}}$ is a power of $\zeta_{n}$.

In view of Remark \ref{remexact}, from Theorem \ref{Ryb2} (applied to $G$ and to $H$), as $\tilde{H}:=\mu_n(GK^\times)=\mu_n(HK^\times)$ we have 
$$\frac{[K(G):K]}
{|GK^{\times}:K^{\times}|}=\delta \cdot \frac{[K(\tilde{H}):K]}{|\tilde{H}:\mu_n(K^{\times})|}=\frac{[K(H):K]}
{|HK^{\times}:K^{\times}|}\,.$$
\end{proof}

\begin{lem}\label{lembetter2}
With the notation of Theorem \ref{better2}, we let $m$ be the largest positive integer such that $\zeta_{2^m}\in H$. If $m=1$ then we have 
$\frac{[K(H):K]}{|HK^\times:K^\times|}=1$, while if $m\geq 2$ then we have 
$$\frac{[K(H):K]}{|HK^\times:K^\times|}=\begin{cases} 
2^{2-m} & \text{if $w>f$ or if $w=f$ and $-\xi_{2^{w+1}}^n\notin G^nK^{\times n}$}\\
2^{1-f} & \text{if $w=f$ and $-\xi_{2^{w+1}}^n\in G^nK^{\times n}$}\\
2^{2-\min(w',m)} & \text{if $w<f$, $\delta=1$}\\
2^{1-\min(w,m)} & \text{if $w<f$, $\delta=1/2$}\end{cases}$$
where, if $w$ is finite, $w'$ is the largest positive integer such that $\zeta_{2^{w'}}\in K(\zeta_4)^\times$.
In the last case we have $m=\max(w, \overline{m})$ where $\overline{m}$ is the largest integer such that $\zeta_{2^{\overline{m}}}\in \mu_n(GK^\times)$.
\end{lem}

We observe the following: in characteristic $0$, we have $w'=w$ or $w'=w+1$ and the latter case holds if and only if $\ch(K)=0$ and $K\cap \Q(\zeta_{{2^\infty}})$ is not totally real; in odd characteristic $p\equiv 3 \bmod 4$, $w'$ is the $2$-adic valuation of $p^2-1$. 

\begin{proof}
If $m=1$, then  $[K(H):K]=|HK^\times:K^\times|$ by Theorem \ref{Kneserthm} and Remark \ref{Knesercondi} so suppose that $m\geq 2$.

We remark that $m\leq w$ if $w>f$ (because  $m\leq f+1$) or if $w=f$ and $-\xi_{2^{w+1}}^n\notin G^nK^{\times n}$. In these cases, we have $[K(H):K]=2$ and $|HK^\times:K^\times|=|H:\mu_{2n}(K^\times)|=2^{m-1}$.

If $w=f$ and $-\xi_{2^{w+1}}^n\in G^nK^{\times n}$, then $K(H)=K(\zeta_4)$.
We conclude because $(1+\zeta_{2^w})^2=\zeta_{2^w}(2+\xi_{2^w})$ and $\zeta_{2^w}$ are in the same class modulo $K^\times$ and hence   
$$|HK^\times: K^\times|=2|\langle \zeta_n\rangle K^\times: K^\times|=n\,.$$

Finally suppose that $w<f$. 
Since $1+\zeta_{2^w}\in K(\zeta_4)^\times$, we have 
$$[K(H):K]=2^{1+\max(m-w',0)}\,.$$

If $\delta=1$, we may conclude because we have $|HK^\times:K^\times|=2^{m-1}$.
If $\delta=1/2$ (in particular, $\ch(K)=0$ and $K\cap \Q(\zeta_{{2^\infty}})$ is  totally real), recall that $(1+\zeta_{2^w})^2\in \zeta_{2^w} K^\times\setminus K^\times $. By Lemma \ref{orders} we know that $(1+\zeta_{2^w})^n=\xi_{2^{w+1}}^{n}\notin K^{\times n}$ so there is no integer $i$ such that 
$(1+\zeta_{2^w})\zeta_{n}^i\in K^\times$.
We deduce that 
$$|HK^\times: K^\times|=2 |\langle \zeta_{2^w}\rangle \mu_n(GK^\times)K^\times: K^\times|\,.$$
To conclude that $|HK^\times: K^\times|=2^m$
we prove that $m=\max(w,\overline{m})$. 
For $\overline{m}\geq w$, $H$ is contained in $K(\zeta_{2^{\overline{m}}})^\times$ and we conclude because this group does not contain $\zeta_{2^{\overline{m}+1}}$. %check-ff 
For $\overline{m}<w$, $H$ is contained in $K(\zeta_{4})^\times$ and we conclude because $\zeta_{2^{w+1}}\notin K(\zeta_{4})^\times$.
\end{proof}

Remark that $K(\zeta_4)=1K+\zeta_4 K$. If $w>f$ or if $w=f$ and $-\xi_{2^{w+1}}^n\notin G^n K^{\times n}$, we only have entanglement if $m\geq 3$. Since $m \leq f+1$ and $m\leq f$ for $w=f$, the $K$-linear relations for $K(G)$ (beyond those stemming from the group $GK^\times/K^\times$) are generated by those expressing 
$$\zeta_{2^3},\ldots, \zeta_{2^m}\in 1 K+\zeta_4K\,.$$
If $w=f$ and $-\xi_{2n}^n\in G^n K^{\times n}$, there is also an additional entanglement (as there is the loss of a factor $2$ in the degree $[K(G):K]$) which is due to $1+\zeta_{2^w}\in GK^\times \cap K(\zeta_4)^\times$, and it is expressed by the $K$-linear relation
\begin{equation}\label{eq:specialent}
1+\zeta_{2^w}\in 1 K+ \zeta_4 K\,.    
\end{equation}
Finally, suppose that $w<f$. If $\delta=1$, then the entanglement is similarly due to
\begin{equation}\label{eq:normalent}
\zeta_{2^3},\ldots \zeta_{2^{m'}}\in 1K +\zeta_4 K
\end{equation}
where $m'$ is the largest positive integer less or equal to $m$ such that $\zeta_{2^{m'}}\in K(\zeta_4)^\times$.
If $\delta=1/2$, the  entanglement is similarly explained by \eqref{eq:normalent} and \eqref{eq:specialent}.


\section{The general case}

Let $n$ be a positive integer coprime to the characteristic of $K$. If $n=1$ then we have 
$$[K(G):K]=|GK^\times:K^\times|=1$$
so we suppose that $n\geq 2$ and write 
$n=\prod_{p} p^{v_p}$ for the prime factorization of $n$, where $p$ varies among the prime divisors of $n$. Let $z$ be the product of the odd primes $p$ such that $\zeta_p\notin K^\times$ and $\zeta_p\in GK^\times$.
We set $n_p:=p^{v_p}$ and $G_p=G^{n/n_p}$. In this way, $n_p$ is the smallest positive integer such that $G_p^{n_p}\in K^\times$. Since $n_p$ is a prime power, we may apply the results in the previous sections to study $G_p$.

\begin{rem}\label{split}
We clearly have
$$|GK^\times:K^\times|=\prod_p |G_pK^\times:K^\times|$$
where $p$ varies among the prime divisors of $n$, and the same holds if we replace $K$ by a finite extension. We also have
$$[K(G):K(\zeta_z)]=\prod_p [K(\zeta_z,G_p):K(\zeta_z)]$$
because by Corollary \ref{corbetterp} (for $p$ odd), by Theorem \ref{Kneserthm} (for $p=2$ and $\zeta_4\in K^\times$ or $4\nmid n$) and by 
Theorem \ref{better2} and Lemma \ref{lembetter2} (in the remaining case) the factors on the right hand side are a power of $p$ and hence the fields $K(\zeta_z,G_p)$, whose compositum is $K(G)$, are linearly disjoint over $K(\zeta_z)$.
\end{rem}

The following result is \cite[Theorem 2]{Schinzel-ab}:

\begin{thm}[Schinzel's theorem on abelian radical extensions]\label{Schinzel-abelian}
Let $n\geq 1$ be not divisible by $\ch(K)$. 
If $a\in K^\times$, the extension $K(\zeta_{n}, \sqrt[n]{a})/K$ is abelian if and only if $a^{m}=b^{n}$ holds for some $b\in K^{\times}$ and for some $m\mid n$ such that $\zeta_{m}\in K$.
\end{thm}

\begin{proof}[Proof of Theorem \ref{mainthm} if $n$ is odd]
By Remark \ref{split} we can write
$$\frac{[K(G):K]}{|GK^\times:K^\times|}=[K(\zeta_z):K]\cdot \prod_{p\mid n}
    \frac{[K(\zeta_z,G_p):K(\zeta_z)]}{|G_pK^\times:K^\times|}\,.
    $$

We have $G_pK^\times \cap K(\zeta_z)^\times \subseteq \mu_{p^{v_p}}(G_p K^\times)$ because $\zeta_p\notin K^\times$ and the extension $K(\zeta_z)/K$ is abelian (we apply Theorem \ref{Schinzel-abelian}).
    
By Theorem \ref{Kneserthm} (in view of Remark \ref{Knesercondi}) we then have 
$$[K(\zeta_z,G_p):K(\zeta_z)]=|G_pK(\zeta_z)^\times:K(\zeta_z)^\times|=\frac{|G_pK^\times:K^\times|}{|\mu_{p^{v_p}}(G_p K^\times)\cap K(\zeta_z)^\times:\mu_{p^{v_p}}( K^\times)|}\,.
    $$
We may then conclude remarking that 
$$|\mu_{n}(G K^\times)\cap K(\zeta_z)^\times:\mu_{n}(K^\times)|=\prod_{p\mid n}|\mu_{p^{v_p}}(G_p K^\times)\cap K(\zeta_z)^\times: \mu_{p^{v_p}}( K^\times)|\,.$$
\end{proof}

\begin{defi}\label{Delta}
We set $\Delta=0$ if $\zeta_4\in K^\times$ or $4\nmid n$. In the remaining case, we let 
$H'$ be the group $H$ from Theorem \ref{better2} and Lemma \ref{lembetter2} for $G^{n'}$ over $K(\zeta_z)$ and set 
\begin{equation}\label{magicDelta}
2^{-\Delta}:=\frac{[K(\zeta_z, H'):K(\zeta_z)]}{|H'K(\zeta_z)^\times:K(\zeta_z)^\times|}\,.
\end{equation}
\end{defi}

\begin{proof}[Proof of Theorem \ref{mainthm} if $n$ is even]
Call $G_{P}=\prod_{p\mid n, p\neq 2} G_p$. Remarking that $\zeta_z\in G_P$, we can write     $$\frac{[K(G):K]}{|GK^\times:K^\times|}=[K(\zeta_z):K]\cdot 
\frac{[K(G_{P}):K(\zeta_z)]}{|G_{P}K^\times:K^\times|} \cdot
    \frac{[K(\zeta_z,G_2):K(\zeta_z)]}{|G_2 K^\times:K^\times|}
    $$
By the odd case of Theorem \ref{mainthm} we have
$$\frac{[K(G_{P}):K]}{|G_{P}K^\times:K^\times|}=\frac{[K(\zeta_z):K]}{|\mu_{n/n_2}(G_P K^\times)\cap K(\zeta_{z})^\times: \mu_{n/n_2}( K^\times)|}\,.$$
Since $\mu_{n/n_2}(G K^\times)=\mu_{n/n_2}(G_P K^\times)$ and $\mu_{2n_2}(G K^\times)=\mu_{2n_2}(G_2 K^\times)$ and $GK^\times \cap \sqrt{K^\times}=G_2K^\times \cap \sqrt{K^\times}$ we are left to prove that
$$\frac{[K(\zeta_z,G_2):K(\zeta_z)]}{|G_2 K^\times:K^\times|}=\frac{2^{-\Delta}}{|\mu_{2n_2}(G_2 K^\times)(G_2K^\times\cap\sqrt{K^\times})\cap K(\zeta_z)^\times: K^\times|}\,.$$
As $K(\zeta_z)/K$ is abelian, by Theorem \ref{Schinzel-abelian}
we have   
$$G_2 K^\times\cap K(\zeta_z)^\times=\mu_{2n_2}(G_2 K^\times)(G_2K^\times\cap\sqrt{K^\times})\cap K(\zeta_z)^\times$$
so it suffices to show that 
$$\frac{[K(\zeta_z,G_2):K(\zeta_z)]}{|G_2 K(\zeta_z)^\times:K(\zeta_z)^\times|}=2^{-\Delta}\,,$$
which is a consequence of Theorem \ref{better2} and \eqref{magicDelta} (or of Theorem \ref{Kneserthm} if $\zeta_4\in K^\times$ or $4\nmid n$).
\end{proof}


\begin{proof}[Proof of Theorem \ref{divisibility}]
Equivalently, we prove that   
$[K(\zeta_z, G):K(\zeta_z)]$ divides 
$\frac{1}{z} \cdot |GK^{\times}:K^{\times}|$.
Letting $p$ be a prime number, by Remark \ref{split} we have 
$$[K(\zeta_z, G):K(\zeta_z)]=\prod_{p\mid n} [K(\zeta_z, G_p):K(\zeta_z)]$$
and 
$$\frac{1}{z} \cdot |GK^{\times}:K^{\times}|=\prod_{p\mid z}\frac{1}{p} \cdot |G_pK^{\times}:K^{\times}| \prod_{p\mid n, p\nmid z} |G_pK^{\times}:K^{\times}|\,.$$
For $p\mid z$ the degree 
$[K(\zeta_z, G_p):K(\zeta_z)]$ divides 
$\frac{1}{p} \cdot |GK^{\times}:K^{\times}|$ by Corollary \ref{corbetterp}. If $p\neq 2$ and  $p\nmid z$, or if $p=2$ and $\zeta_4\in K^\times$ or $4\nmid n$ we have 
$$[K(\zeta_z, G_p):K(\zeta_z)]=|G_pK(\zeta_z)^{\times}:K(\zeta_z)^{\times}|$$ by Theorem \ref{Kneserthm} (in view of Remark \ref{Knesercondi}) and this index divides $|G_pK^{\times}:K^{\times}|$. For $p=2$, $\zeta_4\notin K^\times$ and $4\mid n$ the degree $[K(\zeta_z, G_2):K(\zeta_z)]$ divides $|G_2 K(\zeta_z)^{\times}:K(\zeta_z)^{\times}|$ by Theorem \ref{better2} and \eqref{magicDelta}.
\end{proof}

We set $\mu_\infty=\cup_{m\geq 1}\mu_m$. We conclude by proving a result that shows the eventual maximal growth of certain radical extensions:

\begin{thm}\label{eventual}
For every positive integer $N$ let $R_N$ be a subgroup of $\overline{K}^\times$ such that the index $|R_NK^\times:K^\times|$ divides $N^c$ for some constant $c$, $R_1\in K^\times$ and such that $R_{N}^M=R_{N/M}$ holds for every $M\mid N$. Suppose that there are only finitely many primes $p$ such that $\zeta_p\notin K^\times$ and $\zeta_p\in R_N K^\times$ for some $N$, and call $z$ their product. Moreover, suppose that 
$$|\mu_{\infty}(K(\zeta_{4z})^\times):\mu_{\infty}(K^\times)|$$
is finite.
Then there exists a positive integer $N_0$ such that 
$$\frac{[K(R_N):K]}{|R_NK^\times:K^\times|}=\frac{[K(R_{\gcd(N,N_0)}):K]}{|R_{\gcd(N,N_0)}K^\times:K^\times|}\,.$$
\end{thm}
\begin{proof}
Let $N_0$ be a number that is divisible by $4z$ and with the property that for every $N$ the group $\mu_{N}(R_N K^\times)\cap K(\zeta_{z})^\times$
is a subgroup of $\mu_{N_0}(R_{N_0} K^\times)$. 
Thus removing from $N$ the prime factors coprime to $N_0$ does not affect $\mu_{N}(R_N K^\times)\cap K(\zeta_{z})^\times$. Moreover, if $p$ is any prime number, we have  
$\mu_{p^{v_p(N)}}(R_N K^\times)=\mu_{p^{v_p(N)}}(R_{p^{v_p(N)}} K^\times)$. Combining these two observations we obtain  
$$\mu_{N}(R_N K^\times)\cap K(\zeta_{z})^\times=\mu_{\gcd(N,N_0)}(R_{\gcd(N,N_0)} K^\times)\cap K(\zeta_{z})^\times\,.$$
If $N$ is odd, we may conclude by Theorem \ref{mainthm}. So suppose that $N$ is even. Since  $R_N\cap\sqrt{K^\times}=R_{2^{v_2(N)}}\cap\sqrt{K^\times}$ and because of the bound on $|R_NK^\times:K^\times|$ we may define $N_0$ (such that $v_2(N_0)$ is large enough) so that 
$R_NK^\times\cap\sqrt{K^\times}=R_{\gcd(N,N_0)}K^\times\cap\sqrt{K^\times}$. 
Similarly, we may define $N_0$ such that the group
$$\mu_{2^{v_2(N)+1}}(R_N K^\times)(R_N K^\times\cap\sqrt{K^\times})\cap K(\zeta_z)^\times$$
does not change by replacing $N$ by $\gcd(N,N_0)$ (because the squares of its elements are in $\mu_{2^{v_2(N)}}(R_N K^\times)\cap K(\zeta_z)^\times$ which stabilizes when $v_2(N)$ is large enough). We may then conclude by Theorem \ref{mainthm} because, considering Definition \ref{Delta}, we may define $N_0$ such that $v_2(N_0)>w'$ (or we have $w'=\infty$) and such that $1+\zeta_{2^{w'}}$ is contained in $R_{2^{v_2(N_0)}}$ if it is contained in $R_{2^{v}}$ for some positive integer $v$.
\end{proof}

The following result is the reformulation in our setting of \cite[Theorem 1]{MAMA}:

\begin{thm}\label{MAMA}
Let $K$ be a number field, fix a finitely generated subgroup $\Gamma$ of $K^\times$ and for every positive integer $N$ let $R_N=\sqrt[N]{\Gamma}$. Then there exists a positive integer $N_0$ such that 
$$\frac{[K(R_N):K]}{|R_NK^\times:K^\times|}=\frac{[K(R_{\gcd(N,N_0)}):K]}{|R_{\gcd(N,N_0)}K^\times:K^\times|}\cdot \prod_{p\mid N, p\nmid N_0, \zeta_p \notin K^\times} \frac{p-1}{p}\,.$$
\end{thm}
\begin{proof}
There is an odd squarefree integer $Z$ such that for all primes $p\nmid Z$ we have $[K(\zeta_p):K]=p-1$.
Additionally, we can choose $Z$ such that for any $N\geq 1$ the extensions $K(R_{p^{v_p(N)}}, \zeta_Z)/K(\zeta_Z)$ are linearly disjoint for every prime number $p$. Thus 
for any odd squarefree integer $Z'$ that is a multiple of $Z$ and for every positive integer $N$ we have 
$$\mu_{2^{v_2(N)+1}}(R_N K^\times)(R_N K^\times\cap\sqrt{K^\times})\cap K(\zeta_{Z'})^\times \subseteq R_{2^{v_2(N)}} K^\times\cap K(\zeta_{Z'})^\times \subseteq  K(\zeta_{Z})^\times\,.$$
By Lemma \ref{lembetter2} we may choose the $2$-adic valuation of $N_0$ to be large enough such that 
$$\frac{[K(R_{2^{v_2(N)}}, \zeta_Z):K(\zeta_Z)]}{|R_{2^{v_2(N)}}K(\zeta_Z)^\times:K^\times|}=\frac{[K(R_{2^{\min (v_2(N), v_2(N_0))}}, \zeta_Z):K(\zeta_Z)]}{|R_{2^{\min (v_2(N), v_2(N_0))}}K(\zeta_Z)^\times:K^\times|}\,.$$
Then, following the proof of Theorem \ref{mainthm}, we are left to control those $N$ which divide a power of $Z$, and for them we can find a suitable $N_0$ following the proof of Theorem \ref{eventual}.
\end{proof}

\begin{thebibliography}{0}

\bibitem{Albu} \textsc{Albu,\@ T.}, \emph{Cogalois theory}, Pure and Applied Mathematics 252, Marcel Dekker, New York, 2003. 

\bibitem{BarreraVelez} \textsc{Barrera Mora,\@ F. and Vélez,\@ W.~Y.} \emph{Some results on radical extensions}, J. Algebra \textbf{162} (1993) no.~2, 295--301. 

\bibitem{EntRad} \textsc{Chan,\@ C.W., Pajaziti,\@ A., Perucca,\@ A.  and Perissinotto,\@ F.} \emph{The entanglement of radicals}, arXiv preprint \url{https://arxiv.org/abs/2508.19211}.


\bibitem{Koch1} \textsc{Halter-Koch,\@ F.}, \emph{Eine Galoiskorrespondenz für Radikalerweiterungen (A Galois correspondence for radical extensions)}, J.Algebra \textbf{63} (1980), 318--330.

\bibitem{Koch2} \textsc{Halter-Koch,\@ F.}, \emph{Über Radikalerweiterungen (On radical extensions)}, Acta Arith. \textbf{36} (1980), 43--58. 

\bibitem{Kneser}
\textsc{Kneser,\@ M.}, \emph{Lineare Abhängigkeit von Wurzeln (Linear dependence of roots)}, 
Acta Arith. \textbf{26} (1975), 307--308. 

\bibitem{LangAlgebra}
\textsc{Lang,\@ S.} \emph{Algebra}, Graduate Texts in Mathematics 211, Springer-Verlag, New York, 2002.

\bibitem{Lenstra} 
\textsc{Lenstra,\@ H.\@ W.\@ jr.}, \emph{Entangled radicals}, Colloquium Lectures, AMS 112th Annual Meeting, San Antonio, January 12--15, 2006, available at \url{https://www.math.leidenuniv.nl/~hwl/papers/rad.pdf}.

\bibitem{LMS} 
\textsc{Lenstra,\@ H.\@ W.\@ jr., Stevenhagen,\@ P. and Moree,\@ P.}, \emph{Character sums for primitive root densities.}
Math. Proc. Cambridge Philos. Soc. \textbf{157} (2014), no. 3, 489--511.

\bibitem{Palenstijn} \textsc{Palenstijn,\@ W.\@ J.}, \emph{Radicals in arithmetic}, PhD thesis, University of Leiden (2014), available at \url{https://openaccess.leidenuniv.nl/handle/1887/25833}.

\bibitem{MAMA} \textsc{Perucca,\@ A., Sgobba,\@ P. and Tronto,\@ S.}, \emph{Kummer theory for number fields via entanglement groups}, Manuscripta Math., \textbf{169} (2022), no.~1-2, 251--270.

\bibitem{Rybowicz}
\textsc{Rybowicz,\@ M.}, \emph{On the normalization of numbers and functions
defined by radicals}, J. Symbolic Comput., 
\textbf{35} (2003), 651--672.

\bibitem{Schinzel-ab}
\textsc{Schinzel,\@ A.}, \emph{Abelian binomials, power residues and exponential congruences}, Acta Arith. \textbf{32} (1977) no.~3, 245--274.
Addendum, ibid. \textbf{36} (1980), 101--104. See also Andrzej Schinzel Selecta Vol.II, European Mathematical Society, Z{\"u}rich, 2007, 939--970.

\end{thebibliography}

\end{document}

